\section{Backward induction and decentralised provision}\label{sec:equilibrium}
We first characterise downstream continuations, then test the four independent stage-1 entry incentives. This order matters: the sixteen first-stage profiles do not constitute sixteen exogenous payoffs.

\subsection{Private fee and participation choices}
Fix participation in the other region and define $U_i=u_i+a_i^m m_j>0$, $T_i=\Delta_i+b_i^m q_j>0$. For $s_i=P$, fees and selection satisfy
\begin{equation}
 \ell_i=[(f_i-U_i)/T_i]_0^1,\qquad
 \mathcal R_i(f_i)=N_i f_i(1-\ell_i).
\end{equation}
The global fee optimum on the nonnegative half-line, including the full-participation kink, is
\begin{equation}
 (f_i^*,\ell_i^*)=
 \begin{cases}
 ((U_i+T_i)/2,\ (T_i-U_i)/(2T_i)),&0<U_i<T_i,\\[2mm]
 (U_i,\ 0),&U_i\geq T_i.
 \end{cases}\label{eq:monopolyfee}
\end{equation}
In particular, one must not extend the interior optimum through the kink. If $s_i=B$, free participation is universal and the premium upgrade demand is $N_i(1-f_i/\Delta_i)$ for $f_i\in[0,\Delta_i]$. Hence
\begin{equation}
 f_i^B=\Delta_i/2,\quad D_i^B=1/2,\quad
 \mathcal R_i^B=N_i\Delta_i/4.
 \label{eq:coexist}
\end{equation}

When $(s_A,s_B)=(P,P)$, neither region's participation moments may be held constant while optimising its private fee. A candidate price pair must jointly solve
\begin{equation}
 \ell_i=\left[\frac{f_i-u_i-a_i^m(1-\ell_j)}
 {\Delta_i+b_i^m(1-\ell_j^2)/2}\right]_0^1,\quad i=A,B,
 \label{eq:pp-fixedpoint}
\end{equation}
alongside each provider's global nonnegative-fee best response to $f_j$, allowing both boundary and interior fee branches. On the interior branch the stationary equations may be written in terms of the coupled participation system; stationarity alone is \emph{not} a globally valid fee-equilibrium certificate. We use the documented Stage~4/4A sufficient conditions and numerical best-response diagnostics for the admissible baseline, retaining the possibility of multiple continuations or unresolved cases outside that domain.

\subsection{First-stage equilibrium and planner evaluation}
For each of the sixteen entry configurations, let $\mathcal E(s)$ denote the set of admissible downstream equilibria. Where a continuation is selected, assign the four stage-1 payoffs from \eqref{eq:Wi} and \eqref{eq:privateprofit}. A candidate pure first-stage equilibrium must make every binary entry flip weakly unprofitable:
\begin{align}
 W_i(s)&\ge W_i(s_i^{G\text{-flip}},s_j)
       &&\text{for each }G_i,\label{eq:govBR}\\
 \pi_i^P(s)&\ge \pi_i^P(s_i^{P\text{-flip}},s_j)
       &&\text{for each }P_i.\label{eq:privBR}
\end{align}
The four states in one region are $0\leftrightarrow G$ and $P\leftrightarrow B$ for a government flip, and $0\leftrightarrow P$ and $G\leftrightarrow B$ for a private flip. Thus there are four independent binary incentive inequalities at each state; a government's entry comparison does \emph{not} subsume the private provider's response. Equilibrium ties may generate multiplicity. The restricted planner in \eqref{eq:planner} compares all sixteen total-welfare values \emph{under the same allowed continuation rule}. Without an admissible off-path equilibrium one cannot declare a global first-stage SPNE.

The exhaustive baseline implementation and its numerical limits are documented in the Stage~9 reproduction package. In particular, checks of a coupled $(P,P)$ fee best reply on a finite grid are supplemental evidence, not general uniqueness or a completed Lean proof of the full game.
